\chapter{Conclusiones}
\label{cap:conclusiones}

Este trabajo de tesis tuvo como objetivo principal el estudio de la dinámica de los campos escalares durante la época de \textit{Reheating} y la consecuente producción de un Fondo Estocástico de Ondas Gravitacionales (SGWB). A partir del análisis teórico y la implementación de simulaciones numéricas, se han obtenido diversas conclusiones fundamentales que conectan la cosmología primordial con la astrofísica observacional.

En primer lugar, respecto a la dinámica lineal y la resonancia paramétrica, se desarrolló y validó un marco computacional propio a orden lineal para estudiar la evolución del campo inflatón y la subsecuente creación explosiva de partículas durante la etapa de \textit{Preheating}. Esta implementación permitió comprender fenomenológicamente la estructura de las bandas de inestabilidad y los regímenes de resonancia dictados por el análisis de Floquet. Apoyándose en estas herramientas, se analizó el efecto directo de dichas fluctuaciones lineales sobre el tensor de energía-momento. Se corroboró, tanto analítica como numéricamente, que el crecimiento exponencial de los modos escalares actúa como una fuente eficaz para la amplificación del espectro de potencias de las perturbaciones tensoriales.

Para abordar la dinámica en el régimen no perturbativo, se implementó exitosamente el integrador numérico \CosmoLattice. El uso de redes espaciales tridimensionales permitió superar las limitaciones de la teoría lineal, capturando la evolución completa del sistema. Se lograron simular con éxito los efectos de retroacción (\textit{backreaction}), la dispersión no lineal de modos (\textit{rescattering}) y el vaciamiento total del condensado del inflatón, procesos que dictaminan la saturación y la estabilización definitiva de los espectros de ondas gravitacionales.

Mediante estas herramientas, se llevó a cabo una profunda exploración paramétrica, caracterizando el espacio de soluciones para distintas familias de potenciales. Específicamente, se contrastaron modelos de tipo monomial ($m^2\phi^2$, $\lambda\phi^4$) frente a modelos con \textit{plateau} fenomenológicamente favorecidos, como los atractores $\alpha$ (T-models) y el modelo de Starobinsky. Este barrido demostró de forma concluyente cómo la pendiente del pozo de potencial y la escala de acoplamiento modifican drásticamente la forma del espectro, la frecuencia pico y la amplitud general de la energía tensorial inyectada.

Finalmente, para establecer la sensibilidad observacional del modelo, se reescalaron los espectros teóricos proyectándolos al \textit{strain} característico ($h_c$) observable en la actualidad. Al confrontarlos con las curvas de sensibilidad de arreglos de púlsares (PTA), interferómetros láser y cavidades resonantes electromagnéticas (EMRC), se concluyó que estos últimos experimentos de ultra-alta frecuencia representan la vía más prometedora para la detección de este tipo de señales. Además, se demostró que el paso al espacio del \textit{strain} introduce una ruptura de degeneración para los modelos de tipo materia ($\omega_{eff}=0$): mientras que la densidad de energía $\Omega_{GW}$ de estos modelos resulta virtualmente idéntica a tiempos largos, la dependencia del \textit{strain} con la inversa de la frecuencia magnifica su separación. Esto evidencia que el SGWB posee el poder resolutivo necesario para discriminar empíricamente entre los distintos modelos inflacionarios subyacentes.

\section*{Perspectivas futuras}

Como extensión de este trabajo, resulta de particular interés abordar el estudio de escenarios más complejos que incorporen campos vectoriales de gauge (U(1) o SU(2)) en las simulaciones de \CosmoLattice, con el fin de evaluar cómo las interacciones de norma modifican la eficiencia en la transferencia de energía y la consecuente firma gravitacional. Asimismo, sería valioso ampliar el análisis espectral considerando familias de potenciales polinomiales, tal como se ha propuesto recientemente en la literatura para dinámicas de tipo \textit{hilltop} \cite{Das_2025}. Finalmente, otra dirección prometedora consistiría en enmarcar este estudio dentro de teorías de gravedad modificada, analizando la dinámica y la producción de ondas gravitacionales en escenarios donde el acoplamiento de los campos de materia con la gravedad es no mínimo.